% ===========================================================================
\chapter{Rescaling Predicted Latents to a Reference Norm Removes Most of a State Probe's Error}
\label{chap:bounded}
\margintoc
% ===========================================================================

% ─────────────────────────────────────────────────────────────────────────────
% PUBLICATION HEADER
% ─────────────────────────────────────────────────────────────────────────────

\begin{publicationbasis}{In preparation}
Sager, P.\,J., Taoudi-Benchekroun, Y., Grewe, B.\,F., \& Stadelmann, T.
(2026). Correcting state-probe evaluation of world models for norm drift.
\cite{sager_bounded_2026}
\end{publicationbasis}

% ─────────────────────────────────────────────────────────────────────────────
\contributionstatement
\label{sec:bounded_contribution}
% ─────────────────────────────────────────────────────────────────────────────

The thesis author is the first author of this work. The paper's own author
contributions section records the split, and this statement follows it.

The thesis author conceived the study, designed the method, wrote the
software, carried out the analysis and the investigation, produced the
figures and wrote the first draft. Yassine Taoudi-Benchekroun contributed to
the method and revised the manuscript. Benjamin Grewe and Thilo Stadelmann
supervised the work and revised the manuscript.

This chapter is the second rewrite of a study that began as a positive
result. The original work proposed a spectral constraint on a latent
transition operator and reported large reductions in rollout error. Those
reductions did not survive the controls the thesis author added later, and
the paper was rewritten twice: first as a study of what the constraint really
does to the measurement, and then as the correction to the measurement
protocol that this chapter reproduces. The decision to publish the correction
in place of the original claim is the thesis author's.

\section{Chapter Introduction}

The two chapters before this one each measured a trained network and drew a
conclusion from what the measurement said. Each conclusion rests on the
readout that produced it, and this chapter takes up a case where the readout
the field uses reports something other than what it names.

The setting is a latent world model, which encodes an observation, predicts
forward in its own latent space, and never reconstructs an observation. There
is nothing directly interpretable in that space, so a probe stands in for one.
A regressor is fitted from latent space to physical state on encoded
observations and then turned loose on the latents the model predicts. The quantity the
probe is supposed to report is how well the model predicts the state.

Where the training loss leaves the norm of a latent free, most of what the
probe reports is the drift in that norm. This is the same shape of problem as
the one the previous chapters have to answer about their own measure, and
\cref{chap:discussion} turns the standard set here back on the participation
ratio. The contribution to the thesis argument is therefore twofold. The
chapter takes up a property that training leaves open and that something else
then settles, which is the thesis's subject in a third substrate. It also
supplies the one standard in this thesis that applies to more than one
chapter.

This chapter has no circuit in it and computes no participation ratio. It is
the methodological chapter and is introduced as such. The manuscript it
reproduces is not yet submitted, as the publication box above states.

% ---------------------------------------------------------------------------
% The published paper, included verbatim. Do not edit imported_papers/.
% ---------------------------------------------------------------------------

\begin{importedpaper}{bounded}

% The paper puts its overview figure in the introduction, which this chapter
% drops, and then refers back to it from the results. Reproducing the float
% here, inside the block and under the paper's own label, is what keeps that
% reference resolving. check_imports.py holds it to the original.
%
% [htbp] and not [t], for the reason recorded in sections/3-1_acu_survey.tex:
% a float contributed before any of the paper's own is still pending when the
% environment's group closes, and LaTeX reports Float(s) lost.
\begin{figure}[htbp]
\centering
\includegraphics[width=\linewidth]{figures/fig_intro.pdf}
\caption[Norm drift inflates the probe's error, and re-projection removes it]{\textbf{Norm drift inflates the probe's error, and re-projection removes it.} \textbf{(A)} The prediction is moved along its existing direction to a reference radius estimated from observation encodings. \textbf{(B)} The same TCRL pendulum predictions are evaluated before and after this change. Errors are divided by the error of always predicting the mean state. Persistence predicts no change from the initial state. Shading marks the horizon supervised during training. Diamonds mark the start-encoding probe, fitted from the initial encoding to the state $T$ steps ahead. Longer horizons use separate rollouts under the same policy.}
\label{fig:intro}
\end{figure}

\importpart{03_materials-and-methods}
\importpart{04_results}
\end{importedpaper}

\section{Chapter Summary}

Fit a regressor from latent space to physical state on encoded observations,
apply it to predicted latents, and it responds to the size of what it is given
as well as to the direction, and the training loss of a
decoder free world model need not fix that norm. A cosine loss compares
directions and leaves the norm free. During autoregressive rollout the norm of
a predicted latent then drifts while its direction stays accurate, and the
probe reports the drift.

The size of the effect is large enough to decide a comparison. In synthetic
latents that hold the state in their direction, a norm error of $5\%$ per
prediction step multiplies the norm by ten within $47$ steps, and the probe's
error reaches $27$ times that of always predicting the mean state. When the norm varies
independently of the true error, the probe orders only about $61\%$ of pairs
of latents correctly, where chance is $50\%$. A worse model that drifts less
can therefore rank above a better one, and the chapter shows it happening.
Against models
built to have a known accuracy, the ordinary probe ranks the worse model first
in all $120$ pairs, and re-projection recovers $88$ of those rankings with no
false confirmation.

The same distortion appears in published models. For TCRL on pendulum,
rescaling the predictions to the norm of the encodings removes $97\%$ of the
probe's error at the end of the training horizon, because the directions were
accurate all along and only the norms were wrong. Across cosine trained models
whose dynamics do not constrain the norm, the median removed is $99\%$ twenty
steps ahead. Of $39$ surveyed evaluations of predicted latents, one reports
prediction norms.

The correction is re-projection. Each prediction is rescaled to a reference
norm taken from the encodings, its direction is kept, and the same probe is
applied again. It needs no retraining and can be added to any state probe
evaluation. Four checks come with it, each asking whether the re-projected
error reports the model's prediction or an artefact of the rescaling.

Run on the paper's own trained world models, those checks fail. None of the
rankings between the four dynamics modules holds, because the rollouts do no
better than a probe fitted from the initial encoding, or they land away from
the directions the encoder produces. That is a negative result about the
authors' own models, and it is the strongest evidence in the chapter that the
checks have teeth.

Where the training design already keeps the norm in check, the correction
changes nothing and should not be applied blindly. Squared error training
penalises a wrong norm within the supervised horizon, and there re-projection
raises the error for most module and task combinations, because the norm is
then carrying the model's uncertainty about an uncertain future and removing
it discards information. An output constraint that sets the norm of every
prediction from its direction leaves predictions unchanged under re-projection
by construction. Beyond the supervised horizon the drift returns even under
squared error. The chapter asks that an evaluation establish which quantity
its readout responds to before its ranking is read. Removing norms is a
separate question and the chapter takes no position on it.

\Cref{chap:cna} takes the last of the properties this thesis treats as open
and asks what happens if it is closed by construction instead.

